\section{Discussion}
\label{sec:discussion}

\paragraph{Sources of error.}
All results are obtained on people absent from training, a setting that \cref{sec:ablations} shows
to be considerably harder than a random split suggests. The remaining errors have three main
sources, which call for different remedies: signs that are genuinely ambiguous within a closed
vocabulary (near-synonyms, pointing and directional signs); landmark detection failures, which
dominate the hardest clips; and variation between signers, which exceeds the difference between
the two architectures. More signers and better hand tracking are therefore likely to help more
than a larger classifier.

\paragraph{Absence of transfer from ASL.}
We expected an ASL encoder to help, at least with few LSFB examples, since both languages build
signs from handshapes, movements and locations. The two corpora, however, differ in more than the
language. ASL clips are citation forms recorded one at a time on a smartphone; LSFB clips are
extracted from spontaneous conversations filmed in a studio, are less than half as long (median
\numLsfbDurMedian{}\,s against \numAslDurMedian{}\,s) and are co-articulated with neighbouring signs.
Hands are missing in \numAslNoHandPct{}\,\% of the ASL frames but in \numLsfbNoHandPct{}\,\% of the LSFB
frames, and the landmarks come from two
generations of MediaPipe. The pre-trained encoder is therefore adapted to a different input
distribution, and even a few hundred labelled clips of the target corpus provide more relevant
information than this prior. This explanation is consistent with our measurements but not
demonstrated: separating the effect of the language from that of the recording conditions would
require a second ASL corpus of conversational signs, or an LSFB corpus of citation forms.
Self-supervised pre-training on unlabelled, in-domain continuous signing, for instance LSFB-CONT
\citep{fink2021lsfb} or MEDIAPI-SKEL \citep{bull2020mediapi} for French Sign Language, appears to be
a more promising direction than supervised cross-lingual pre-training.

\paragraph{Methodological pitfalls.}
Several of our observations are inexpensive to check yet easily overlooked. (i) Image-space
translation and scaling augmentations are cancelled by body-centred normalisation and must be
applied after it. (ii) Normalised landmark coordinates are anisotropic whenever pixels or frames
are not square, which distorts the body frame across datasets. (iii) Early stopping on a plateau
while the learning rate is still high can make an ablation appear much more harmful than it is.
(iv) Rare feature outliers, harmless in floating point, can break post-training quantisation, and
quantisation does not necessarily make a small model faster on CPU. (v) A preprocessing function
implemented in two languages must be tested for parity, which costs little once fixtures are
generated from the reference implementation.

\paragraph{Limitations.}
(i) The ASL test set contains only four signers, so per-signer estimates are uncertain, and all
LSFB clips come from a single corpus recorded in one studio setting. (ii) Both vocabularies are
small and fixed: an out-of-vocabulary sign is always mapped to the closest known sign, and the demo
has no rejection class. (iii) Our ASL accuracy remains well below that of the best competition
solutions \citep{kaggle_islr}, which relied on stronger regularisation and augmentation, much
longer schedules and ensembles of several models, and were trained on all 21 participants, whereas
our models see 14 of them and are tested on a different set. We did not measure how much each of
these factors contributes. A longer schedule alone did not help our model, and the ablations used a
single seed, \numAblEarlyStopped{} of them ending before the full schedule; differences of about one
point between ablations should therefore not be over-interpreted. (iv) Because every clip is
stretched to 64 frames, velocities are expressed per resampled frame: their scale depends on the
duration of the clip and on the frame rate of the source (30 or 50\,fps), which may partly explain
why padding performed better in the ablations. Resampling to a common frame rate and padding would
remove this dependence. (v) The lip contour is the only non-manual information used; eyebrows, gaze and head
movements, which carry grammatical information, are ignored. (vi) The demo segments signs from hand
presence and motion, which suits signs produced one at a time with a pause but not natural,
continuous signing. (vii) No Deaf signer took part in the evaluation of the demo; its usefulness
for the community remains to be established with its members.

\paragraph{Ethical considerations.}
Both datasets were released for research by their creators, under the Kaggle competition rules and
the CC BY 4.0 licence (LSFB-ISOL). ASL data are never redistributed; the demo ships a few LSFB clips
with attribution. The demo processes video locally and stores nothing. Face and body landmarks can
nonetheless identify a person, so any extension that records them should obtain explicit consent.
We describe the system as \emph{isolated sign recognition}: presenting such a classifier as a
translator would mislead users about its capabilities and could cause harm where communication
matters. A model card summarising intended uses, out-of-scope uses and known biases accompanies
the code \citep{mitchell2019modelcards}.
